\head{Grades K--1: towers and one empty dot}
Below, $0$ is the empty dot and other numbers are tower heights, listed along a path from the left. Plain-words key for the whole packet: \emph{the tallest tower always stands next to the empty dot.} Mirror images (the same towers read from the other end) count as different ways unless the children decide otherwise; either choice is fine if they say which they use.

\sub{Problem 1, page 1: path of 3 dots, towers 1 and 2}
Exactly $4$ ways: $1$-$0$-$2$, $2$-$0$-$1$, $0$-$2$-$1$, $1$-$2$-$0$ (two ways and their mirror images). Complete: a link of $2$ needs the $2$-tower beside the empty dot, and the $1$-tower then goes on the far side of either one. Ask ``How do you know you have found them all?'' Hint: ``Which tower has to stand next to the empty dot?''

\sub{Problem 2, page 2: path of 4 dots, towers 1, 2, 3}
Exactly $4$ ways: $0$-$3$-$1$-$2$, $3$-$0$-$2$-$1$ and their mirror images $2$-$1$-$3$-$0$, $1$-$2$-$0$-$3$. Complete: the $3$-tower stands beside the empty dot. If those two take the middle pair of dots, the end towers make links $1,1$ or $2,2$; so they sit at one end, and only one order of the other two towers works each way round. Hint: ``Can the $3$ and the empty dot both be in the middle?''

\sub{Problems 3 and 4, page 3: path of 5 dots and the fork}
Path: $8$ ways, for example $0$-$4$-$1$-$3$-$2$ (links $4,3,2,1$): empty, tallest, shortest, next tallest, \dots. Fork: $12$ ways, for example empty dot on the left end, $4$ on the dot where the path branches, $2$ then $3$ along the rest of the row, $1$ below the branch (links $4,2,1,3$). Hint for both: ``Where can the $4$ go?'' Extension: find a second way.

\sub{Problem 5, page 4: star with 3 links}
Finish the small pictures with middle $0$ (empty) and middle $3$; the other three pieces go round the outside in any order. Cross out middles $1$ and $2$. A child's reason: the $3$-link needs the $3$-tower and the empty dot together, and every link touches the middle, so one of them must be the middle. With $1$ in the middle the outside pieces $0,2,3$ make links $1,1,2$; with $2$, links $2,1,1$. Hint: ``Which two have to be joined to make the $3$?''

\sub{Problem 6, page 5: triangle, towers 1, 2, 3}
Circle $1$ and $2$. Empty, $1$, $3$ gives links $1,2,3$ with the $2$ left over; empty, $2$, $3$ gives $2,1,3$ with the $1$ left over. The $3$ is never left over, since nothing else makes a $3$-link. Every two dots of a triangle are joined, so once the three pieces are right, any placement works. Hint: ``Which two pieces could make the $3$? Can the $3$-tower stay in the tray?''

\sub{Problem 7, page 6: square, towers 1, 2, 3, 4}
Circle $1$ and $3$. Going round: empty, $2$, $1$, $4$ (links $2,1,3,4$) leaves the $3$; empty, $4$, $2$, $3$ (links $4,2,1,3$) leaves the $1$. The $4$ must be used. If the $2$ is left over, the empty dot and the $4$ are joined, the $1$ and $3$ are joined to each other, and the two remaining links measure $1,1$ or $3,3$. ``How do you know the $2$ cannot be left over?'' is a good question for a child who circled only $1$ and $3$. Hint: ``Where must the $4$-tower stand?''

\sub{Problem 8, page 7: star with 4 links}
The $4$-tower, or no tower: the middle can also be the empty dot. Either way the other four pieces go outside in any order ($24$ ways each). The reason is Problem~5's. Because the question says ``which towers,'' children may miss the empty middle; ask ``Could the middle be the empty dot?''

\head{Grades 2--3: cards, loops and a ruler}
\sub{Problem 1, page 1: paths of 4, 5 and 6 dots}
From the left: $0,3,1,2$; \ $0,4,1,3,2$; \ $0,5,1,4,2,3$. There are $4$, $8$ and $24$ labellings in all. Many children find the zigzag (smallest, largest, next smallest, next largest). Hint: ``Where can the biggest card go?''

\sub{Problem 2, page 2: path of 11 dots}
$0,10,1,9,2,8,3,7,4,6,5$ along the snake from the top left ($648$ labellings in all). Hint: ``What did your paths in Problem 1 have in common?''

\sub{Problem 3, page 3: two trees}
Fork: $0$ on the left end, $4$ on the branch dot, $2$ then $3$ along the row, $1$ below (links $4,2,1,3$; $12$ in all). Two joined middles with two outer dots each: one middle $0$ with outer dots $1$ and $4$, the other middle $5$ with outer dots $2$ and $3$ (links $1,4,5,3,2$; $48$ in all). Hint: ``Which card must be next to the $0$?''

\sub{Problem 4, page 4: middle of the 5-link star}
Only $0$ or $5$, each with the other cards outside in any order. A link of $5$ can only be made by $0$ and $5$, and every link touches the middle, so one of them is the middle. For the other cards the explanation is the same each time: with $2$ in the middle, a $5$-link would need a $7$, which is not a card. Hint: ``Which two cards can make the link of $5$?''

\sub{Problem 5, page 5: loops of 3, 4, 5 and 6 links}
Label the loops of $3$ and $4$; cross out $5$ and $6$. Going round: $0,1,3$ or $0,2,3$ (the $2$ or the $1$ left over); $0,2,1,4$ or $0,4,2,3$ (the $3$ or the $1$ left over). The loops of $5$ and $6$ have no labelling (a computer tried every arrangement); children cross them out after many failures, which is evidence, and Problem 7 asks for the reason. If a child is stuck on the $5$-loop, ask ``Which card are you leaving out? Have you tried leaving out a different one?'' and later ``Do you think it can be done?''

\sub{Problem 6, page 6: loop of 7 links}
For example $0,7,1,6,2,4,3$ going round: the path zigzag with the $5$ left out ($168$ labellings in all). Only about $1$ in $240$ arrangements of the cards works, so random trying is slow. Hint: ``Which two cards must be neighbours?'' Only as a last resort, since it hands over the construction: ``Could you start like a path, $0,7,1,6,\dots$, and leave out one card?''

\sub{Problem 7, page 6: why the crossed-out loops fail}
The climbing count from the overview. A third-grade route: walk a finger round the loop; each link is a step up or a step down by its number, and you come back to the same card, so the ups add up to the same as the downs. All the link numbers together then make an even total, but $1+2+3+4+5=15$ and $1+2+\dots+6=21$ are odd. On a labelled square the climbs and falls are $1,4$ and $2,3$ in some order, and $1+4=2+3$. A held prompt: ``You came back to where you started. What do you know about the ups and the downs?'' Follow-ups: ``Does this show that the $3$- and $4$-loops can be labelled?'' (no; it only fails to rule them out) and ``Which loop is the next one it rules out?'' ($9$ and $10$). This is the stretch problem of the packet; a construction plus ``we tried a lot'' is a fine place to stop.

\sub{Problem 8, page 7: four counters on the strip from 0 to 6}
Counters at $0,1,4,6$ or $0,2,5,6$ (mirror images). Complete: a distance of $5$ needs a counter at $1$ (with $6$) or at $5$ (with $0$). With $1$, the distances so far are $1,5,6$, and the last counter must make $2,3,4$: only $4$ does ($2$ repeats $1$, $3$ makes two 3s, $5$ repeats $1$). With $5$, the mirror image gives $2$. Hint: ``Which squares could make the distance $5$?'' Extension: put $0,1,4,6$ on four dots that are all joined to each other, and check the six links (Grades 4--5 Problem 8).

\head{Grades 4--5: rules, proofs, rulers and trees}
\sub{Problem 1, page 1}
Path of 6 from the left: $0,5,1,4,2,3$ ($24$ labellings). Two joined middles: as in Grades 2--3 Problem 3. Hint: ``Where can the $5$ go?''

\sub{Problem 2, page 2: any path}
$0,10,1,9,\dots,5$ along the snake. The rule $0,m,1,m-1,2,\dots$ uses each label from $0$ to $m$ once, and its gaps are $m,m-1,\dots,1$ because the low labels climb by one while the high labels fall by one. A complete explanation says both why the link numbers are all different and why they run from $1$ to $m$. Hint: ``Read the link numbers along your path. What happens from one link to the next?'' Other rules exist ($648$ labellings for $11$ dots). Extension: does the idea work for a star, or for a path with extra dots hanging off it (the caterpillar of Problem 9)?

\sub{Problem 3, page 3, and Problem 4, page 4: loops and the proof}
Answers as Grades 2--3 Problem 5, and the climbing count of the overview, in either form. Held prompt for Problem 4: ``If you walk round the loop and come back to the start, what do you know about all the ups and all the downs?'' A child who argues with odd and even numbers has an equally good proof. Ask whether the argument shows that the $3$- and $4$-loops \emph{can} be labelled; it does not.

\sub{Problem 5, page 5: loops of 7 and 8 links}
Both can be labelled. The climbing count allows them ($28$ and $36$ are even), and constructions settle them: $0,7,1,6,2,4,3$ and $0,8,1,7,2,5,3,4$ going round ($168$ and $384$ labellings; about $1$ in $240$ and $1$ in $945$ random arrangements). Hint: ``Which two cards must be neighbours?'' Only as a last resort, since it hands over the construction: ``What happens if you go round the loop with a path labelling? Which link goes wrong, and what if you leave out a different card?''

\sub{Problem 6, page 6: 9, 10, 11, 12, 100 and 101 links}
The proof rules out $9$, $10$ and $101$: $1+\dots+9=45$, $1+\dots+10=55$ and $1+\dots+101=5151$ are odd. For $11$, $12$ and $100$ the totals $66$, $78$ and $5050$ are even, so the proof says nothing, and in fact they can be labelled: $0,11,1,10,2,9,3,7,4,6,5$ and $0,12,1,11,2,10,3,8,4,7,5,6$, and for $100$ the same zigzag skipping $75$. So ``I cannot show it'' is the honest answer for those three. Pairing $1+100$, $2+99$, \dots\ gives $50\times101$ quickly. Hints: ``Which of these totals do you already know?'' ``Is there a quick way to add $1$ to $100$?'' Extension: which numbers of links does the proof rule out? (Those leaving remainder $1$ or $2$ on division by $4$.)

\sub{Problem 7, page 7: rulers}
Four counters: $0,1,4,6$ or $0,2,5,6$, as Grades 2--3 Problem 8. Five counters from $0$ to $10$: impossible. A child's case argument: a distance of $9$ needs a counter at $1$ or $9$; by mirror image say $1$. A distance of $8$ then needs a counter at $8$, since $2$ would make a second $1$ and $9$ a second $9$. Now $0,1,8,10$ give $1,2,7,8,9,10$, and the last counter must make $3,4,5,6$; each empty square repeats a distance (try $2$ to $7$ and $9$). A shortcut: a counter at $x$ between $2$ and $7$ adds distances $x$, $x-1$, $8-x$ and $10-x$, which always total $17$, but $3+4+5+6=18$. The odd-and-even proof of the overview is the adult's version. Hint for five counters: ``With counters on $0$ and $10$, where could a distance of $9$ come from?'' Five counters go on the materials \emph{Strip from 0 to 10} (page 1); the page's $0$--$10$ strips ($1.7$\,cm squares) are too small for counters and serve for writing down tries.

\sub{Problem 8, page 8: four and five dots, all joined}
Four dots: $0,1,4,6$ placed anywhere (all $24$ placements work, since every two dots are joined), or $0,2,5,6$. Five dots: impossible, and Problem 7 is the reason, since the labels would be five counters from $0$ to $10$ with every distance once. A second proof reuses Problem 4: trace the pentagon and then the star, a trip using every link once that comes home, while $1+\dots+10=55$ is odd. Hints: ``How is this like Problem 7?'' ``Can you travel along every link once and come back?''

\sub{Problem 9, page 9: trees}
The spider: middle $0$; one leg $6$ then $1$, another $4$ then $3$, the third $2$ then $5$ (links $6,5$; $4,1$; $2,3$; $60$ labellings in all). The caterpillar: left end $0$, the dot with four neighbours $7$, its upper and lower dots $1$ and $2$, the next dot along $3$ with $6$ below it, the last dot $5$ with $4$ below it (links $7,6,5,4,3,2,1$; $144$ in all). That is the path zigzag run along the spine, and children can turn it into a rule for every caterpillar: $0$ on the end dot of a longest path and the top card on its neighbour; that dot's other neighbours get the next low cards, ending with the next dot along the spine; that dot's new neighbours get the next high cards, again ending on the spine; and so on. Each new link is one shorter than the last. The printed trees' dots are too small for cards: lay cards on the materials \emph{Trees} sheet (page 3), or write numbers on the page. Hint: ``Which card must be next to the $0$?'' The spider is the smallest tree that is not a caterpillar, so the rule does not reach it. ``Can you find a tree that cannot be labelled?'' No one had by the 2013 survey (see the overview), and every tree with up to $10$ dots was checked here. A child who cannot label a tree has a candidate, not a counterexample: the table can try it together, and the adult can ask which trees were hardest and why.

\head{Sources and use}
Alexander Rosa, ``On certain valuations of the vertices of a graph,'' in \emph{Theory of Graphs} (International Symposium, Rome, 1966), Gordon and Breach, 1967, pp.~349--355: the loop theorem. Solomon W. Golomb, ``How to number a graph,'' in \emph{Graph Theory and Computing}, ed. R.~C. Read, Academic Press, 1972: complete networks and rulers. Joseph A. Gallian, ``A Dynamic Survey of Graph Labeling,'' \emph{Electronic Journal of Combinatorics} DS6, 16th ed. (2013): \S2.1 pp.~7 and~9 (trees; the conjecture open; trees up to 35 dots checked), \S2.2 p.~9 (cycles, Rosa), \S2.4 pp.~16 and~18 (complete graphs, Golomb rulers). JRMF, \emph{Graceful Gourds} (web page) and its \emph{Graceful Graphs Festival Guide}, pp.~2--3, a model of a short rule demonstration (see \texttt{plans/atlas/DESIGN.md}). These come from the theme's outline and were not re-read for this guide; check the current survey before saying the conjecture is open. R.~Montgomery, A.~Pokrovskiy and B.~Sudakov, ``A proof of Ringel's conjecture,'' arXiv:2001.02665 (2020), located for this guide; Keevash and Staden's independent proof is cited from memory. Every count, labelling and impossibility above was recomputed from the final pages by \texttt{check\_guide.py} in this guide's source folder; the zigzag loops up to $1000$ links and the caterpillar rule on random caterpillars are computer checks. Library links: Week~10 (walks over every bridge and odd islands), Week~29 (which lengths rods build on a strip) and Week~55 (cards and sums on a strip).
